\documentclass[11pt,a4paper]{amsart}
\usepackage[T1]{fontenc}
\usepackage{lmodern}
\usepackage{amssymb}
\usepackage{microtype}
\usepackage{booktabs}
\usepackage[hyphens]{url}
\usepackage[colorlinks=true,linkcolor=black,citecolor=black,urlcolor=black]{hyperref}
\usepackage[margin=2.5cm]{geometry}

\hypersetup{
  pdftitle={Odd-cycle colourings of abelian Cayley graphs come from characters},
  pdfauthor={Sergi Gal\'an},
}

\newtheorem{theorem}{Theorem}
\newtheorem{proposition}[theorem]{Proposition}
\newtheorem{lemma}[theorem]{Lemma}
\newtheorem{corollary}[theorem]{Corollary}
\theoremstyle{definition}
\newtheorem{question}[theorem]{Question}
\newtheorem{remark}[theorem]{Remark}
\newtheorem{example}[theorem]{Example}

\newcommand{\Q}{\mathbb{Q}}
\newcommand{\R}{\mathbb{R}}
\newcommand{\Z}{\mathbb{Z}}
\newcommand{\Cay}{\operatorname{Cay}}
\newcommand{\RZ}{\R/\Z}

\title{Odd-cycle colourings of abelian Cayley graphs come from characters}
\author{Sergi Gal\'an}
\thanks{These results were obtained with the assistance of an AI system (Claude, by Anthropic), which carried
out the computations and the checks and helped to draft this note. A separate AI agent refereed the draft, with
its own checker and tests, and its findings are incorporated. No mathematician has checked it yet. The author is
not a professional mathematician and submits it for independent verification.}
\date{Draft, 3 October 2026}
\subjclass[2020]{Primary 05C15; Secondary 05C25, 11R11, 52C10}
\keywords{Cayley graph, abelian group, $3$-colouring, circular colouring, distance graph, unit-distance graph,
Hadwiger--Nelson problem}

\begin{document}

\begin{abstract}
We show that a Cayley graph of an abelian group has a homomorphism to the odd cycle $C_{2k+1}$ if and only if some
character maps every generator into the closed arc $[k/(2k+1),(k+1)/(2k+1)]$ of $\RZ$. In particular such a graph
is $3$-colourable if and only if some character maps every generator into $[1/3,2/3]$. The proof averages the
discrete winding number of Krebs and Sankar over the group with an invariant mean. Consequences: Payan's theorem
that a cube-like graph is never $3$-chromatic, and the exponent-$4$ case of the theorem of Krebs and Sankar, follow
at once; a distance graph $G(\Z,D)$ is $3$-colourable if and only if $\kappa(D)\ge1/3$; and for a number field $F$
the plane $F^2$ is $3$-colourable if and only if a character of the group generated by the unit vectors maps them
into $[1/3,2/3]$. It also follows that in every abelian group, every set of Bohr recurrence is a set of
$3$-chromatic recurrence: the case of three colours of a question of Katznelson. For $F=\Q(\sqrt d)$ and finitely
many unit vectors the last condition is a question on a
$4$-dimensional torus, which exact branch-and-bound certificates settle. We obtain $\chi(\Q(\sqrt d)^2)=4$ for
$d=83,107,203$, $4\le\chi(\Q(\sqrt{143})^2)\le5$ and $\chi(\Q(\sqrt{167})^2)\ge4$: five fields for which searches
for finite non-$3$-colourable unit-distance graphs had failed.
\end{abstract}

\maketitle

\section{Introduction}

For an abelian group $\Gamma$ and a subset $S=-S$ of $\Gamma\setminus\{0\}$, the Cayley graph $\Cay(\Gamma,S)$ has
vertex set $\Gamma$ and an edge between $g$ and $g+s$ for $g\in\Gamma$, $s\in S$. A \emph{character} is a
homomorphism $\xi:\Gamma\to\RZ$. For $x\in\RZ$ let $\|x\|$ be its distance to $0$.

\begin{theorem}\label{thm:main}
Let $k\ge1$ and $n=2k+1$. The graph $\Cay(\Gamma,S)$ has a homomorphism to the cycle $C_n$ if and only if there is
a character $\xi$ of $\Gamma$ with $\xi(s)\in[k/n,(k+1)/n]$ for every $s\in S$. In particular, $\Cay(\Gamma,S)$ is
$3$-colourable if and only if some character maps every element of $S$ into $[1/3,2/3]$.
\end{theorem}

A homomorphism to $C_{2k+1}$ is a circular $(2k+1)/k$-colouring. Let $\kappa(S)=\sup_\xi\inf_{s\in S}\|\xi(s)\|$;
the supremum is attained, as the dual group is compact and $\xi\mapsto\inf_s\|\xi(s)\|$ is upper semicontinuous.
The theorem says
\[
\chi_c(\Cay(\Gamma,S))\le 2+\tfrac1k \iff \kappa(S)\ge \tfrac{k}{2k+1}\qquad(k\ge1).
\]
The \emph{if} direction is the classical inequality $\chi_c\le1/\kappa$: for distance graphs it is Zhu's
multiplier method, surveyed by Liu \cite[\S2]{Liu}, and for Cayley graphs of $\Z/p$ it appears in
\cite[Theorem~3.2]{Alon}. The content of the theorem is the converse at the thresholds $2+1/k$, among them
$\chi\le3$.

Steps~1 and~2 of the proof are the discrete winding number of Krebs and Sankar \cite{KS} (their Definition~3.5,
Propositions~3.3 and~3.6, Remark~3.7, and the parity argument in the proof of their Theorem~4.1), which they use to
show that groups of exponent $4$ have no $3$-chromatic Cayley graphs. What is new, as far as we know, is Step~3,
which averages the winding over the group, and the characterisation that results (Section~\ref{sec:lit}).

For $k=1$ the theorem has a reading in terms of recurrence. A set $S\subseteq\Gamma$ is a set of
\emph{$r$-chromatic recurrence} if every $r$-colouring of $\Gamma$ has two points of one colour that differ by an
element of $S$, and a set of \emph{Bohr recurrence} if it meets every Bohr neighbourhood of $0$. Katznelson's
question \cite{Katznelson,Griesmer} asks whether every set of Bohr recurrence is a set of $r$-chromatic recurrence
for every $r$; it is open, and the answer is not known for any countably infinite abelian group \cite{Griesmer}.
For $r=2$ the answer is easily yes. Theorem~\ref{thm:main} gives $r=3$:

\begin{corollary}\label{cor:bohr}
A set $S\subseteq\Gamma\setminus\{0\}$ is a set of $3$-chromatic recurrence if and only if it meets the set
$\{g:\|\xi(g)\|<1/3\}$ for every character $\xi$. In particular, every set of Bohr recurrence is a set of
$3$-chromatic recurrence, and every $3$-colourable abelian Cayley graph has a $3$-colouring whose colour classes are
$\xi^{-1}[0,1/3)$, $\xi^{-1}[1/3,2/3)$ and $\xi^{-1}[2/3,1)$ for a single character $\xi$.
\end{corollary}

\begin{proof}
Apply Theorem~\ref{thm:main} with $k=1$ to $S\cup(-S)$; the colouring is the one in the \emph{if} direction of the
proof.
\end{proof}

We have not found the case $r=3$ stated in \cite{Griesmer,HKM,LWYZ}; we have not been able to read
\cite{Katznelson}.

Section~\ref{sec:proof} proves the theorem. Section~\ref{sec:exp} derives Payan's theorem \cite{Payan} and the
exponent-$4$ case of Krebs and Sankar \cite{KS}. Section~\ref{sec:dist} treats distance graphs.
Section~\ref{sec:planes} applies the theorem to the unit-distance graph of the plane over a number field, the
setting that led to it, and Section~\ref{sec:four} records that nothing of the kind holds for four colours.

\section{Proof of Theorem~\ref{thm:main}}\label{sec:proof}

\emph{If.} Let $c(g)=2\lfloor n\,\xi(g)\rfloor\bmod n$, where $\xi(g)$ is read in $[0,1)$. Along an edge, $n\xi$
moves by an amount in $[k,k+1]$ modulo $n$, so $\lfloor n\xi\rfloor$ moves by $k$ or $k+1$ modulo $n$, and $c$
moves by $2k\equiv-1$ or $2k+2\equiv1$. So $c$ is a homomorphism to $C_n$, whose vertices are $\Z/n$ with
$j\sim j\pm1$.

\emph{Only if.} Let $c:\Gamma\to\Z/n$ be a homomorphism to $C_n$, and let $\Gamma'$ be the subgroup generated by
$S$. A character of $\Gamma'$ extends to $\Gamma$, since $\RZ$ is divisible, so it suffices to find the character
on $\Gamma'$. We work in $\Gamma'$ because $\Cay(\Gamma',S)$ is connected, which Step~2 uses: averaging over all
of $\Gamma$ would fail when $\Cay(\Gamma,S)$ is disconnected, since different components can wind in opposite
directions. For $g\in\Gamma'$ and $s\in S$ let $\sigma(g,s)\in\{1,-1\}$ be the lift of $c(g+s)-c(g)$; then
$\sigma(g+s,-s)=-\sigma(g,s)$.

\smallskip
\emph{Step 1: squares do not wind} \cite[Proposition~3.6]{KS}. For $s,t\in S$ the two integers
$\sigma(g,s)+\sigma(g+s,t)$ and $\sigma(g,t)+\sigma(g+t,s)$ lie in $\{-2,0,2\}$ and are congruent to
$c(g+s+t)-c(g)$ modulo $n$. Two distinct elements of $\{-2,0,2\}$ differ by $2$ or $4$, which $n$ (odd, at
least $3$) does not divide; so they are equal.

\smallskip
\emph{Step 2: the winding of a closed walk} \cite[Definition~3.5, Propositions~3.3 and~3.6,
Remark~3.7]{KS}. Let $Z$ be the abelian group generated by symbols $[s]$, $s\in S$, subject to $[-s]=-[s]$. It is
free on one element of each pair $\{s,-s\}$, except that $2[s]=0$ when $2s=0$; so the parity of the number of
symbols is a well-defined homomorphism $Z\to\Z/2$. Let $R$ be the kernel of $Z\to\Gamma'$, $[s]\mapsto s$.

A closed walk from $g$ with steps $s_1,\dots,s_N$ ($\sum s_j=0$) has the sum
$\sum_j\sigma(g+s_1+\dots+s_{j-1},s_j)\equiv0\pmod n$. By Step~1 this sum does not change when two consecutive
steps are exchanged, nor when a step $s$ immediately followed by $-s$ is deleted. Every walk can be brought to a
normal form by these moves (sort the steps by exchanges, then cancel the pairs $s,-s$), so the sum depends only on
the base point and on the class $r=\sum_j[s_j]\in R$. It does not depend on the base point either: going from
$g+s$ to $g$, doing the walk and coming back gives the same sum, and $\Cay(\Gamma',S)$ is connected. Write the sum
as $n\,w_c(r)$. Then $w_c:R\to\Z$ is a homomorphism (concatenate walks). Since $n\,w_c(r)$ is a sum of $N$ odd
numbers and $n$ is odd, $w_c(r)$ has the parity of $r$.

\smallskip
\emph{Step 3: averaging.} As $\Gamma'$ is abelian, it is amenable: there is a mean $m$ on the bounded functions on
$\Gamma'$ that is invariant under translations. For finite $\Gamma'$ it is the plain average, and for finitely
generated $\Gamma'$ a limit along F\o lner sets. Let
\[
f(s)=m\bigl(g\mapsto\sigma(g,s)\bigr)\in[-1,1].
\]
Then $f(-s)=-f(s)$. For a closed walk $(s_1,\dots,s_N)$ that represents $r\in R$, the function
$g\mapsto\sum_j\sigma(g+s_1+\dots+s_{j-1},s_j)$ is the constant $n\,w_c(r)$ by Step~2. Applying $m$, which is linear
and invariant under translations, gives
\[
\sum_{j=1}^N f(s_j)=n\,w_c(r).
\]

\smallskip
\emph{Step 4: the character.} Put $\xi(s)=\tfrac12\bigl(f(s)/n+1\bigr)\in[k/n,(k+1)/n]$; then
$\xi(-s)=1-\xi(s)$. If $2s=0$, then $\sigma(g+s,s)=-\sigma(g,s)$ for all $g$, so $f(s)=0$ and $\xi(s)=1/2$. Hence
$\xi$ is well defined modulo $1$ on $Z$. For $r\in R$ represented by a walk of length $N$,
$\sum_j\xi(s_j)=\tfrac12\bigl(w_c(r)+N\bigr)$, an integer by the parity in Step~2. So $\xi$ vanishes on $R$ modulo
$1$ and defines a character of $Z/R=\Gamma'$, with $\xi(S)\subseteq[k/n,(k+1)/n]$. \qed

\begin{remark}
For $\Gamma=\Z^2$ with the standard generators this is the height-function picture of $3$-colourings of the
square lattice, in which the height changes by $\pm1$ along edges and by $0$ around squares (compare
\cite{Sheffield}); $f$ is the averaged slope. Theorem~\ref{thm:main} says that the averaged slopes of any colouring
of any abelian Cayley graph form a character as above. The colouring $c$ itself need not come from a character: for
$\Gamma=\Z^2$ most $3$-colourings do not.
\end{remark}

\section{Groups of small exponent}\label{sec:exp}

\begin{corollary}\label{cor:exp}
Let $\Gamma$ be an abelian group of exponent $2$ or $4$. Then no Cayley graph of $\Gamma$ has chromatic number $3$.
\end{corollary}

\begin{proof}
A character of $\Gamma$ takes values in $\tfrac14\Z/\Z$, and the only such value in $[1/3,2/3]$ is $1/2$. So if
$\Cay(\Gamma,S)$ is $3$-colourable, a character maps $S$ to $1/2$, and $g\mapsto\lfloor2\xi(g)\rfloor$ is a proper
$2$-colouring.
\end{proof}

For exponent $2$ this is Payan's theorem on cube-like graphs \cite{Payan}, which Cervantes and Krebs reproved
\cite{CK2}. For exponent $4$ it is a case of the theorem of Krebs and Sankar \cite[Theorem~1.2]{KS}: an abelian
group has a $3$-chromatic Cayley graph if and only if its exponent is not $1$, $2$ or $4$. Their proof uses the
same winding number: for exponent $4$ they show that the discrete fundamental group is torsion, and by their
Theorem~4.1 a non-bipartite graph with a torsion fundamental group is not $3$-colourable. The same argument as for
Corollary~\ref{cor:exp} gives, for instance, that a Cayley graph of an abelian group of exponent $8$ that maps to
$C_5$ is bipartite, since $1/2$ is the only element of $\tfrac18\Z/\Z$ in $[2/5,3/5]$. For $C_3$ the values $3/8$
and $5/8$ are also available, and indeed such groups have $3$-chromatic Cayley graphs \cite{KS}.

\section{Distance graphs}\label{sec:dist}

For a finite set $D$ of positive integers, the distance graph $G(\Z,D)$ is $\Cay(\Z,\pm D)$. Its characters are
$x\mapsto\alpha x$, $\alpha\in\RZ$, and $\kappa(D)=\max_\alpha\min_{d\in D}\|\alpha d\|$ is the parameter of the
lonely runner problem \cite{Liu,PS}.

\begin{corollary}\label{cor:dist}
$G(\Z,D)$ is $3$-colourable if and only if $\kappa(D)\ge1/3$; it maps to $C_{2k+1}$ if and only if
$\kappa(D)\ge k/(2k+1)$.
\end{corollary}

The set of $\alpha$ with $\|\alpha d\|\ge1/3$ for all $d\in D$ is a finite union of closed intervals whose
endpoints are among the points $(j+1/3)/d$ and $(j+2/3)/d$. So it is non-empty if and only if it contains one of
these points, and the test is finite. We compared it, for all $1\,747$ sets $D\subseteq[1,24]$ with $|D|=3$ and
$\gcd(D)=1$, with $3$-colourability of the segment $[0,12\max D+59]$ decided by a SAT solver. They agree, and the
$74$ sets that need four colours are exactly those of Zhu's classification \cite{Zhu2002} of $3$-element sets with
$\chi(G(\Z,D))=4$, namely $\{1,2,3m\}$ and $\{a,b,a+b\}$ with $a\not\equiv b\pmod3$ (after division by the gcd).
Liu's survey asks whether $\chi_c(G(\Z,D))<1/\kappa(D)$ can happen when $|D|=3$ \cite[Problem~3]{Liu}. By
Corollary~\ref{cor:dist}, if one of the two numbers is at most $3$, then both equal $2$ or both lie in the same
interval $(2+1/(k+1),2+1/k]$.

\section{Planes over number fields}\label{sec:planes}

Let $F$ be a field of characteristic $0$ with $i\notin F$, and $L=F(i)$. Identify $F^2$ with $L$ by
$(x,y)\mapsto x+iy$; the unit vectors are $T=\{z\in L: z\bar z=1\}$, and the unit-distance graph of $F^2$ is
$\Cay(L,T)$, a disjoint union of copies of $\Cay(A,T)$ with $A=\Z[T]$.

\begin{corollary}\label{cor:planes}
$\chi(F^2)\le3$ if and only if some character of $A$ maps every unit vector into $[1/3,2/3]$. If for a finite set
$U\subset T$ no character of the group $\Z U$ maps $U$ into $[1/3,2/3]$, then $\chi(F^2)\ge4$, and some finite
unit-distance graph over $F$ is not $3$-colourable.
\end{corollary}

The last assertion is the theorem of de Bruijn and Erd\H{o}s \cite{dBE}. If $F$ is a number field, then $\Z U$ is
a finitely generated subgroup of the $\Q$-vector space $L$, hence free of rank at most $[L:\Q]$; for a quadratic
field $F$ the test for a finite $U$ is a question on a torus of dimension at most $4$. A locally constant
$3$-colouring at a place of $F$ \cite{Fischer1990,Moorhouse,LG} is a colouring of a finite quotient of $A$, and
Theorem~\ref{thm:main} applied to that quotient gives such a character.

\subsection{Real quadratic fields}
For $F=\Q(\sqrt d)$ with $d>1$ squarefree, $\chi(F^2)=2$ if $d\not\equiv3\pmod4$
\cite{Johnson1987,Fischer1990,Moorhouse}, $\chi(F^2)\le3$ if moreover $d\not\equiv2\pmod3$
\cite{Fischer1990,Moorhouse}, and $\chi(F^2)\le4$ if $d\equiv3\pmod8$ \cite[Theorem~10]{Fischer1990} or
$d\equiv0,1,2,4\pmod7$ \cite{Moorhouse}. In \cite{QP} we found finite non-$3$-colourable unit-distance graphs over
$\Q(\sqrt d)$ for $27$ values $d\equiv11\pmod{12}$ below $1000$, but searches for such graphs failed for
$d=83,107,143,167,203$, the values that the local--global question of \cite{LG} predicted to need four colours and
that remained open below $210$.

Fix a denominator $D$ and let $U_D$ be the unit vectors $\bigl((a+b\sqrt d)/D,(c+e\sqrt d)/D\bigr)$, $a,b,c,e\in\Z$,
one of each pair $\pm u$. A character with $\xi(U_D)\subseteq[1/3,2/3]$ gives $f\in[1/3,2/3]^{U_D}$ ($f_u$ the lift
of $\xi(u)$ in $[1/3,2/3]$) with $\langle r,f\rangle\in\Z$ for every integer relation $r$ among the vectors of
$U_D$. Using an LLL-reduced basis of the relations, we rule this out by branch and bound over the integers
$\langle r,f\rangle$: each leaf carries a rational vector $y$ (a Farkas certificate) such that
$\langle y,\cdot\rangle$ applied to the equations fixed at that leaf cannot hold on the box $[1/3,2/3]^{U_D}$. A
short program checks a certificate in exact integer and rational arithmetic: the vectors have length $1$
($a^2+db^2+c^2+de^2=D^2$, $ab+ce=0$), the relations hold, every branch covers every integer value in its range, and
every Farkas vector is valid. Any set of valid relations gives necessary conditions, so infeasibility is a proof
whether or not the relations form a basis.

\begin{theorem}\label{thm:fields}
$\chi(\Q(\sqrt d)^2)=4$ for $d=83,107,203$; $4\le\chi(\Q(\sqrt{143})^2)\le5$; and $\chi(\Q(\sqrt{167})^2)\ge4$.
\end{theorem}

\begin{proof}
The lower bounds are Corollary~\ref{cor:planes} with the certificates of Table~\ref{tab:certs} \cite{Repo}. The
upper bound $4$ holds because $83\equiv107\equiv203\equiv3\pmod8$ \cite[Theorem~10]{Fischer1990}; see also
\cite[Theorem~2.2]{ACLMSS}. For $d=143$ the place above $11$ is ramified in $\Q(\sqrt{143})$ with residue field
$\mathbb{F}_{11}$, and $i$ is not in its completion; unit vectors are integral there and reduce to the circle
$x^2+y^2=1$ of $\mathbb{F}_{11}^2$, whose graph is $5$-colourable \cite{QP}, which colours each coset of the
valuation ring.
\end{proof}

\begin{table}[h]
\centering
\begin{tabular}{rrrrrr}
\toprule
$d$ & $D$ & vectors & relations & nodes & leaves\\
\midrule
83 & 510 & 54 & 50 & 6\,886 & 4\,664\\
107 & 1170 & 90 & 86 & 12\,785 & 9\,618\\
143 & 1740 & 54 & 50 & 3\,861 & 2\,710\\
167 & 1560 & 54 & 50 & 8\,552 & 6\,094\\
203 & 1530 & 78 & 74 & 1\,072 & 732\\
\bottomrule
\end{tabular}
\caption{Certificates that no character of $\Z U_D$ maps $U_D$ into $[1/3,2/3]$ (vectors counted up to sign).
In each case the vectors span a group of rank $4$ and the relations listed generate all relations.}
\label{tab:certs}
\end{table}

With the values of \cite{QP}, $\chi(\Q(\sqrt d)^2)$ is now known for every squarefree $d<143$ except $d=47$.
$\Q(\sqrt{167})$ is the first admissible field of \cite{LG}: no place has a locally constant $4$-colouring, so
$\chi(\Q(\sqrt{167})^2)\ge4$ is also the three-colour case of the local--global question there. For $d=203$ the
result shows that no level of the place above $7$ is $3$-colourable, which \cite{LG} had checked only up to level
$3$. For $d=83$ and $D=510$, colouring-guided growth had stopped at $6\,008$ points because a $3$-colouring of the
ball of radius $2$ extended to all $215\,478$ candidate points, so the obstruction does not lie in a small ball. A
subset of $28$ of the $54$ vectors, found by greedy deletion, already suffices for $d=83$. For $\Q(\sqrt{11})$, a
subset of $16$ vectors with $D=30$ suffices, in agreement with the $76$-vertex graph of \cite{QP}.

The lower bounds rest on Theorem~\ref{thm:main}, whose averaging step is not constructive, and on the
certificates; we give no explicit finite graph for the five fields. We expect such graphs to be large, which would
explain why the searches failed, but this is speculation: Step~3 says nothing about their size. A quantitative
version of Step~3 over F\o lner boxes, or an explicit certified graph, would remove the dependence on the infinite
averaging.

All five denominators are even, as they must be. For odd $D$ every vector of $U_D$ has $a+b+c+e$ odd, because
$a^2+db^2+c^2+de^2=D^2$ and $d$ is odd; so $x\mapsto\frac12(a+b+c+e)$, for $Dx=(a+b\sqrt d)+i(c+e\sqrt d)$, is a
character of $\Z U_D$ with the value $1/2$ on $U_D$, and $\Cay(\Z U_D,U_D)$ is bipartite. The infeasible
denominators we found are all multiples of $6$, but some multiples of $6$ are feasible (for instance $D=1170$ for
$d=203$). As controls, the test is feasible for $40$ pairs $(d,D)$ with $d\in\{2,3,5,6,7,15,19,31,39,43\}$, whose
planes are known to be $3$-colourable, and for the even denominators $(d,D)=(107,1450)$, $(143,1752)$,
$(167,1820)$ and $(203,1590)$, whose unit vectors also span a group of rank $4$.

\begin{question}
Is $\chi(\Q(\sqrt d)^2)\ge4$ for every squarefree $d\equiv11\pmod{12}$? In terms of Corollary~\ref{cor:planes} and
\cite{LG}: for real quadratic fields, does a character that maps all unit vectors into $[1/3,2/3]$ exist only when
a place gives a locally constant $3$-colouring?
\end{question}

\section{Four colours}\label{sec:four}

Step~1 fails for $K_4$, and so does the theorem: $K_4=\Cay((\Z/2)^2,(\Z/2)^2\setminus\{0\})$ is $4$-colourable, but
every character of $(\Z/2)^2$ vanishes on some nonzero element, so none maps the three generators into $[1/4,3/4]$.
The inequality $\chi_c\le1/\kappa$ holds for all $k$, so a character with $\xi(S)\subseteq[1/4,3/4]$ still gives a
$4$-colouring. For the $108$ unit vectors of $\Q(\sqrt{47})$ with denominator $240$, which underlie the search for
a five-chromatic triangle-free unit-distance graph in \cite{LG}, a floating-point computation finds none; we have
not certified this.

\section{Literature}\label{sec:lit}

The closest work is that of Krebs and Sankar \cite{KS}, whose winding number is Steps~1 and~2; they use it for
torsion fundamental groups, without characters or averaging. The \emph{if} direction of Theorem~\ref{thm:main} is
classical \cite{Liu,Alon,GKM}. The following are consistent with Theorem~\ref{thm:main} and do not imply it:
Cervantes and Krebs \cite{CK1,CK2,CK3} (in \cite{CK3}, $3$-colourability for small dimension and rank is
characterised by two forbidden subgraphs, diamond lanyards and $C_{13}(1,5)$, and $C_{13}(1,5)$ indeed has no
suitable character); Krebs and Leyva \cite{KL}; the work on Katznelson's question on Bohr recurrence
\cite{Griesmer,HKM,LWYZ}, where we found no statement for three colours (we have not seen \cite{Katznelson}
itself); and Perarnau and Serra \cite[\S3.5]{PS}, which has $\chi_f\le\chi_c\le1/\kappa$. We did
not find Theorem~\ref{thm:main} in these works, in the literature on distance graphs and the lonely runner problem
through Liu's survey \cite{Liu}, or in web searches on characters, circular colourings and odd-cycle
homomorphisms of Cayley graphs. The proof uses only classical tools, so an earlier occurrence is possible.

By Theorem~\ref{thm:main}, whether a finite abelian Cayley graph, given by $\Gamma$ and $S$, maps to $C_{2k+1}$ can
be decided by checking the $|\Gamma|$ characters, with $O(|\Gamma|\,|S|)$ operations: polynomial in the size of
the graph.

\begin{thebibliography}{99}
\bibitem{Alon} N. Alon, \emph{The chromatic number of random Cayley graphs}, European J. Combin. \textbf{34}
(2013).
\bibitem{ACLMSS} M. Axenovich, J. Choi, M. Lastrina, T. McKay, J. Smith and B. Stanton, \emph{On the chromatic
number of subsets of the Euclidean plane}, Graphs Combin. \textbf{30} (2014).
\bibitem{CK1} J. Cervantes and M. Krebs, \emph{Chromatic numbers of Cayley graphs of abelian groups: a matrix
method}, Linear Algebra Appl. \textbf{676} (2023), 277--295; arXiv:2303.06262.
\bibitem{CK2} J. Cervantes and M. Krebs, \emph{An alternate proof of Payan's theorem on cubelike graphs},
arXiv:2303.06267 (2023).
\bibitem{CK3} J. Cervantes and M. Krebs, \emph{Chromatic numbers of Cayley graphs of abelian groups: cases of
small dimension and rank}, arXiv:2303.06272 (2023).
\bibitem{dBE} N. G. de Bruijn and P. Erd\H{o}s, \emph{A colour problem for infinite graphs and a problem in the
theory of relations}, Nederl. Akad. Wetensch. Proc. Ser. A \textbf{54} = Indag. Math. \textbf{13} (1951),
369--373.
\bibitem{Fischer1990} K. G. Fischer, \emph{Additive $K$-colorable extensions of the rational plane}, Discrete
Math. \textbf{82} (1990), 181--195.
\bibitem{Repo} S. Gal\'an, \emph{The Hadwiger--Nelson problem over number fields}, code, data and formal
proofs, 2026. Archived at Zenodo, \href{https://doi.org/10.5281/zenodo.22976635}{doi:10.5281/zenodo.22976635}
(all versions); source at \url{https://github.com/decalion89/chromatic-number-of-the-plane}; the certificates and
checker cited here are in \texttt{data/quadratic\_planes/winding/}, and the account is
\texttt{notes/winding\_lemma.md}.
\bibitem{LG} S. Gal\'an, \emph{Two-colourable planes over number fields}, draft, 2026, in \cite{Repo}.
\bibitem{QP} S. Gal\'an, \emph{Real quadratic planes that need four colours}, draft, 2026, in \cite{Repo}.
\bibitem{GKM} I. Garc\'ia-Marco, K. Knauer and G. Menara, \emph{Coloring semiminimal Cayley graphs},
arXiv:2607.26942 (2026), Lemma~2.2.
\bibitem{Griesmer} J. T. Griesmer, \emph{Special cases and equivalent forms of Katznelson's problem on
recurrence}, arXiv:2108.02190 (2021).
\bibitem{HKM} B. Host, B. Kra and A. Maass, \emph{Variations on topological recurrence}, Monatsh. Math.
\textbf{179} (2016), 57--89.
\bibitem{Johnson1987} P. D. Johnson Jr., \emph{Two-colorings of real quadratic extensions of $\Q^2$ that forbid
many distances}, Congr. Numer. \textbf{60} (1987), 51--58 (not seen by us).
\bibitem{Katznelson} Y. Katznelson, \emph{Chromatic numbers of Cayley graphs on $\Z$ and recurrence},
Combinatorica \textbf{21} (2001), 211--219 (not seen by us).
\bibitem{KL} M. Krebs and A. Leyva, \emph{Chromatic numbers of rank-two Abelian Cayley graphs}, arXiv:2511.03028
(2025).
\bibitem{KS} M. Krebs and M. Sankar, \emph{Abelian groups without $3$-chromatic Cayley graphs}, J. Combin. Theory
Ser. B (2026); arXiv:2410.11028.
\bibitem{Liu} D. D.-F. Liu, \emph{From rainbow to the lonely runner: a survey on coloring parameters of distance
graphs}, Taiwanese J. Math. \textbf{12} (2008), 851--871.
\bibitem{LWYZ} H. Liu, Z. Wu, N. Yang and S. Zhang, \emph{Chromatic thresholds for linear equations and
recurrence}, arXiv:2603.05490 (2026).
\bibitem{Moorhouse} G. E. Moorhouse, \emph{On the chromatic numbers of planes}, preliminary draft, revised
3 March 2010, \url{https://www.ericmoorhouse.org/pub/chromatic.pdf}.
\bibitem{Payan} C. Payan, \emph{On the chromatic number of cube-like graphs}, Discrete Math. \textbf{103} (1992),
271--277.
\bibitem{PS} G. Perarnau and O. Serra, \emph{The lonely runner conjecture turns 60}, arXiv:2409.20160 (2024).
\bibitem{Sheffield} S. Sheffield, \emph{Random surfaces}, Ast\'erisque \textbf{304} (2005).
\bibitem{Zhu2002} X. Zhu, \emph{Circular chromatic number of distance graphs with distance sets of cardinality $3$},
J. Graph Theory \textbf{41} (2002), 195--207.
\end{thebibliography}

\end{document}
